\documentclass[identity=acme]{swisstex}

\setrunningtitle{Acme \textbar{} Technischer Bericht}
\hypersetup{pdftitle={Acme Demonstrationsbericht}}

\swissmeta{docid=TR-2026-014, version=1.2, classification=internal}

\begin{document}

\swisscover*{variant=report, kicker=Technischer Bericht,
  title=Acme Demonstrationsbericht,
  subtitle=Referenzdokument der Identitätsschicht,
  foot=Acme AG · acme.example}

\section{Ausgangslage}

Acme AG dokumentiert ihre technischen Systeme mit SwissTeX. Die
Identitätsdatei \swisscode{swissidentity-acme.sty} bindet Farbe, Schrift,
Logo-Dateien und Klassifizierungsvokabular einmalig an, ausschliesslich
über die öffentlichen Setzer der Klasse; dieses Dokument selbst kennt nur
seine eigenen Angaben aus \swisscode{\textbackslash swissmeta}.

\marg{Ein kühles Rot statt des wärmeren Klassenrots}
Der vorliegende Bericht trägt die Kennung TR-2026-014 in Version 1.2 und
ist als intern eingestuft. Beide Angaben erscheinen automatisch im
Seitenfuss und auf dem Umschlag, weil \swisscode{\textbackslash swissmeta}
sie zentral hält und beide Stellen denselben Wert lesen. Die
Klassifizierung folgt dabei dem geordneten Vokabular, das
\swisscode{swissidentity-acme.sty} über
\swisscode{\textbackslash swissclassifications} festlegt; ein Wert
ausserhalb dieses Vokabulars wäre ein Fehler, kein stiller Rückfall auf
eine Vorgabe. Ebenso liest der Umschlag dieselbe Kennung wie die
Fusszeile, ohne dass der Dokumenttext beide Stellen einzeln pflegen
müsste.\footnote{Die Fusszeilenvorlage der Identität kombiniert Kennung,
Version und Datum zu einer Zeile; ein Dokument ohne eigenes
\swisscode{\textbackslash swissmeta} zeigt an dieser Stelle nichts an.}

\begin{swisstable}{Kennzahlen des Berichts}
\begin{tabularx}{\textcolumn}[t]{@{}S{30mm}@{\hspace{4mm}}L@{}}
\toprule
\tabhead{Merkmal} & \tabhead{Wert} \\
\midrule
Kennung & TR-2026-014 \\
\addlinespace[3.5pt]
Version & 1.2 \\
\addlinespace[3.5pt]
Klassifizierung & Intern \\
\addlinespace[3.5pt]
Identität & acme (Referenzidentität) \\
\bottomrule
\end{tabularx}
\end{swisstable}

Nach der Tabelle setzt der Grundtext wieder auf dem Raster auf. Logo und
Kolophon schliessen den Bericht.

\swisslogo[axis=full]{acme-logo.pdf}
\colophon{Gesetzt mit SwissTeX und der Referenzidentität acme. Grundschrift
SwissTeX Grotesk, 9,5/13,5\,pt im Flattersatz, Satzspiegel 105\,mm.}

\end{document}
